\color{unidarkgreen}{\subsection{Газовая защита РПН}}
\color{black}

\begin{enumerate}

\item
Для защиты устройства РПН применяется струйное реле (в составе ГЗ РПН), которое реагирует на ускоренный поток масла, возникающий при повреждениях в баке РПН.

\item
Предусмотрена возможность приема сигналов срабатывания ГЗ РПН и низкой изоляции цепи ГЗ РПН, по разным цепям от двух источников ПДС1 и ПДС2. Для приема сигналов от ПДС1 предусмотрены входные сигналы \textit{«Ср. контакта ГЗ РПН 1 цепь»} и \textit{«Низкая изол. ГЗ РПН 1 цепь»} для каждой фазы РПН, а для приема сигналов от ПДС2 предусмотрены входные сигналы \textit{«Ср. контакта ГЗ РПН 2 цепь»} и \textit{«Низкая изол. ГЗ РПН 2 цепь»} для каждой фазы РПН. Если сигналы срабатывания ГЗ РПН заводятся непосредственно со шкафа ШОМ силового трансформатора, то конфигурируются только входные сигналы \textit{«Ср. контакта ГЗ РПН 1 цепь»} и \textit{«Низкая изол. ГЗ РПН 1 цепь»} для каждой фазы РПН, а действие ГЗ РПН по второй цепи выводится из работы с помощью уставок \textbf{«ГЗ РПН 2 цепь»}.

\item
Ввод в работу ГЗ РПН осуществляется с помощью уставок \textbf{«ГЗ РПН 1 цепь»} и \textbf{«ГЗ РПН 2 цепь»}. Оперативный перевод ГЗ РПН на сигнал производится с помощью команд управления \textit{«Перевод ГЗ РПН на сигнал»}.

\item
Оперативный ввод/вывод ГЗ РПН производится с помощью БУ (блока управления функциями) \textit{«Вывод ГЗ РПН»}.

\item
Непрерывный контроль изоляции цепи струйного реле ГЗ РПН каждой фазы осуществляется с помощью реле контроля изоляции (РКИ). В случае снижения сопротивления изоляции РКИ подает сигнал \textit{«Низкая изол. ГЗ РПН фаза 1 цепь»} для блокировки действия ГЗ РПН. Время срабатывания блокировки от сигналов РКИ задается с помощью уставки \textbf{«Тбл ГЗ РПН»}. Ввод блокировки ГЗ РПН от РКИ осуществляется с помощью уставки \textbf{«Блок ГЗ РПН»}. Для сброса блокировки от РКИ, используется команда управления \textit{«Возврат ГЗ РПН»}.


\item
Алгоритм работы ГЗ РПН приведен на рисунке \ref{pic:GZRPN}. Параметры для настройки ГЗ РПН приведены в таблице \ref{app_set:gzrpn}.

\vspace{2mm}
\begin{figure}[H]
\centering
\includegraphics[width=1\textwidth,height=1\textheight,keepaspectratio]{\fbpath/06.02 ДЗАТ/041.1071 ГЗ РПН/img/GZ_RPN.pdf}
\captionof{figure}{Функциональная схема логики ГЗ РПН}\label{pic:GZRPN}
\end{figure}


\end{enumerate}